\section{Background and Related Work}
\label{sec:background}
In this section, we first clarify key terms used in our scoping review. We then review existing secondary research on human experiences in smart buildings, identifying gaps that motivate our research.

\subsection{Clarifying Relevant Terminology Pertaining to the Scoping Review}
Our scoping review is built around the terms discussed below: 

\textbf{\textit{``Smart'' Buildings:}} Previous works use terms like ``smart buildings'' and``intelligent buildings'', often interchangeably with the term ``built environment'' \cite{buckman2014smart, latifah2020smart}. Irrespective of the term used, the goal is to create a responsive environment where users interact with buildings through automated systems that optimise functions like lighting, heating, and ventilation, using data-driven approaches to adapt to occupants' needs \cite{saputra2023smart}. Smart buildings also emphasise human-centred design, aiming to create environments that foster occupants' satisfaction or agency \cite{mitchell2020no, mitchell2020human}. In our review, we do not distinguish between these terms, collectively referring to them as ``smart buildings'' throughout the paper.

\textbf{\textit{``Interactions'' (with)in Smart Buildings:}} We view interaction in smart buildings as a dynamic relationship in which humans and buildings influence and determine each other's behaviour. We believe this relationship can encompass all forms of interaction described by \citet{hornbaek2017interaction}: the interaction between occupants and the embedded technologies within a building can be understood as a continuous dialogue, a transmission of information, as tool use, as an adaptation towards optimal behaviour, and a controlled move towards a target state, but also as embodiment and experience (and likely other concepts as well). Using ``interaction'' as another search keyword in our review, we do not discriminate against a specific understanding of interaction. With this view of interaction, in which ``[u]sers, with their goals and pursuits, are the ultimate metric of interaction'' \cite{hornbaek2017interaction}, human experiences can similarly be viewed holistically.

\textbf{\textit{``Human Experiences'' in Buildings:}} There are different conceptualisations of experience, e.g. \citet{wright2003making}'s compositional, sensual, emotional, and spatio-temporal ``threads'' or \citet{wright2008empathy}'s lens of empathy. However, we do not employ a specific existing framework of user experience to constrain the type of human experiences that we extract from the literature. Instead, our review seeks to include diverse and emergent aspects of human experiences in smart buildings, without imposing any predefined structure. We broadly consider the experience of occupants in smart buildings as context-dependent, influenced by embodied cognition, emotions, sensory inputs, and ongoing interactions with their environment. 

\textbf{\textit{``Comfort'':}} In this paper, we frequently refer to ``comfort'' as a technical term that has been meticulously studied in architecture and building performance engineering, and which comprises four dimensions: ``thermal'', ``visual'', ``acoustic'', and ``respiratory''. These dimensions have also emerged in HBI research. For example, thermal comfort has examined personalised heating, ventilation, and air conditioning (HVAC) systems~\cite{chhaglani2022flowsense, khare2023thermal}, addressing social dynamics like shared control \cite{von2021towards, mazzone2021socially} and the impact of data displays with environmental information \cite{clear2018thermokiosk}. Visual comfort studies have explored the importance of lighting design \cite{boyce2003human, bian2023simulation, zheng2023daylighting}. 
Research on indoor air quality has focused on raising awareness~\cite{zhong2020hilo} and encouraging behavioural changes, often using playful systems \cite{snow2019performance, gaver2013indoor}, and acoustic comfort studies examined noise preferences and control mechanisms through surveys \cite{marques2020real, passero2012acoustic} and sound awareness systems \cite{wozniak2012maintaining, olry2020prisme}. These dimensions collectively contribute to the understanding of comfort as a key experience in the built environment. For the rest of this paper, when referring to comfort, we refer to these four dimensions. 

\subsection{Secondary Research on Experiences in the Built Environment}
Several works have reviewed literature on human comfort in built environments. For example, \citet{frontczak2011literature} surveyed how lighting and temperature influence human comfort. Similarly, \citet{zhang2012factors} discuss occupant satisfaction with comfort dimensions in the context of green office buildings. Few works explore human experiences beyond comfort in smart buildings through secondary research (i.e. scoping or systematic reviews). For example, a review by \citet{pourzolfaghar2016investigating} identifies design challenges in creating smart environments but only addresses user experience in broad terms. \citet{aliero2022systematic} explore trends in control over automated buildings and smart technologies while considering user experience in the context of privacy in automation systems. Work by \citet{saputra2023smart} highlights the key components that make up smart buildings, emphasising energy management and building information modelling, while viewing occupant comfort in the context of improving quality of life. Similarly, \citet{latifah2020smart} connect smart building features with smart city frameworks to enhance comfort and transparency in meeting digital society's needs \cite{latifah2020smart}.

Other reviews focus on specific environments. For example, \citet{backlund2023building} examined campus buildings to explore the relationship between occupant behaviour, building systems, and energy use, proposing a framework for assessing the impact of occupant behaviour on energy consumption. \citet{ji2023systematic} reviewed sensing technologies across six areas---occupancy status, physiological indicators, building components, environment, consumption, and multi-sensor fusion---to identify concerns related to privacy and data concealment from building occupants. Smart home reviews have focused on how users view benefits arising from technology in these spaces  \cite{marikyan2019systematic}, and the complexity of domestic spaces and opportunities around human-home collaboration \cite{mennicken2014today}. Further secondary research has examined smart technology adoption by ageing populations \cite{ma2022smart} and communication challenges within smart homes for vulnerable groups \cite{jamwal2022smart}. 
To the best of our knowledge, our scoping review is the first of its kind to encompass all types of smart environments and consider the full spectrum of human experiences beyond traditional notions of comfort.